% LSPV-3 REPORT::FULLSKOLEM1::2026-10-03T07:31:00-07:00
\documentclass[11pt]{article}
\usepackage[margin=.85in]{geometry}
\usepackage{amsmath,amssymb}
\begin{document}
\section*{LSPV-3 FULLSKOLEM1: full PH-level reconstruction search collapses under the P=NP counterfactual}
\textbf{STATUS: OPEN.} P versus NP remains open. This report proves a conditional closure theorem. It proves no SAT lower bound, no arbitrary-algorithm necessity theorem, and no nonrelativizing or nonalgebrizing hinge.

\section*{CORPUS REVIEWED}
This run surveyed the current repository/search surfaces and read root main.tex, AGENTS.md, CODEX.md, the repository pnp-research skill, verified MAIN 0058 and newer synthesis/search surfaces, current task references, materially relevant LSPV-1 through LSPV-5 reports, the HEC/CRS-PPC/CCSPC random-challenge branch, governance notes, and reconstruction reports RAMCOMP1, ENVELOPE1, RAMDOM1, CROSSDOM1, TOTALENV1, BOTTLENECK1 together with the red-team GLOBALGAME1 and STATICNP1 controls. The LSPV-2 two-cycle gap is already documented by research/governance/2026-10-03_0621_LSPV2-TWO-CYCLE-WRITE-GAP.tex; no exact missing scientific body is available and none was reconstructed.

The primary-literature refresh checked Viola's probabilistic-time/alternation simulation, van Melkebeek--Mocelin Sdroievski's instance-wise AM derandomization architecture, and Akshay et al.'s QBF Skolem/Herbrand extraction. These sharpen model boundaries but do not supply an unrestricted SAT lower bound.

\section*{Exact reconstruction predicate}
Fix one source-typed SAME\_REALIZATION, including the shared size-normalized \(Q\), Arthur randomness convention, SAT YES/NO semantics, and model-matched polynomial-time predicate \(R\). The reconstruction target is
\[
\Psi(x,Q)\equiv
\exists \sigma\;\forall u\;\exists z\;\forall y\;
R(x,Q,\sigma,u,z,y).
\]
Assume \(\sigma,u,z,y\) all have lengths polynomially bounded in \(N=|(x,Q)|\), and \(R\) is decidable in deterministic polynomial time in the declared model.

\section*{FULLSKOLEM1 theorem}
\textbf{Theorem (PROVED/EXACT, conditional on \(P=NP\)).}
If \(P=NP\), then on every input \((x,Q)\) satisfying \(\Psi(x,Q)\):
\begin{enumerate}
\item a valid outer witness \(\sigma\) can be constructed in deterministic polynomial time; and
\item after fixing that \(\sigma\), there is a deterministic polynomial-time algorithm \(F\) such that for every \(u\),
\[
\forall y\;R(x,Q,\sigma,u,F(x,Q,\sigma,u),y).
\]
\end{enumerate}
Thus the semantic acquisition of the complete polynomial-time Skolem evaluator is polynomial under the contradiction hypothesis whenever the exact reconstruction obligation is exhausted by the displayed polynomial-hierarchy predicate.

\paragraph{Proof.}
Under \(P=NP\), the polynomial hierarchy collapses to \(P\). Hence the language
\[
L_\sigma=\{(x,Q,p):\exists s\ \forall u\ \exists z\ \forall y\
R(x,Q,ps,u,z,y)\},
\]
with \(ps\) restricted to the polynomial outer-witness length, is in \(P\). Starting from a true instance, ordinary prefix self-reduction queries \(L_\sigma\) once per bit and constructs a valid \(\sigma\) in polynomial time.

Fix the resulting \(\sigma\). For a challenge \(u\) and candidate prefix \(p\) of \(z\), define
\[
L_z=\{(x,Q,\sigma,u,p):\exists s\ \forall y\
R(x,Q,\sigma,u,ps,y)\}.
\]
This is a \(\Sigma_2^P\) language before collapse and therefore lies in \(P\) under \(P=NP\). Because the fixed \(\sigma\) satisfies \(\forall u\exists z\forall y R\), every \(u\) has an initial accepting extension. Prefix self-reduction on \(L_z\) constructs \(z\) in polynomial time. The resulting uniform procedure is \(F\). \(\square\)

\section*{What the theorem does and does not eliminate}
FULLSKOLEM1 is an exact quantifier-elimination closure under the contradiction hypothesis, but only qualitatively at polynomial scale. It does not provide the quantitative fixed-decider exponent required by RAMCOMP1, ENVELOPE1, or EXPINV1. In particular, \(P=NP\Rightarrow PH=P\) does not by itself identify the exponent of the resulting \(P\)-algorithm.

The theorem does prove a stronger guard:
\[
\boxed{\text{displayed PH reconstruction semantics alone cannot force superpolynomial evaluator acquisition under }P=NP.}
\]
Therefore any claimed superpolynomial acquisition obstruction must contain an additional requirement not exhausted by \(\Psi\), and that requirement must be shown both necessary for every relevant static SAT decider and model-matched to the original SAT computation.

\section*{Relation to STATICNP1 and the route envelope}
STATICNP1 closes polynomial SAT staticizations of NP-level continuation winning. FULLSKOLEM1 is not a duplicate: it closes the entire displayed
\[
\exists\sigma\forall u\exists z\forall y
\]
reconstruction semantics directly via the full PH collapse and constructs both the outer \(\sigma\) and the uniform \(u\mapsto z\) selector.

Consequently, the ENVELOPE1 evaluator route
\[
F(K)=\max\{f,c_FK\}
\]
cannot use semantic Skolem acquisition itself as a superpolynomial obstruction in a contradiction proof assuming \(P=NP\). A quantitative comparison still requires source-transcribed exponents; FULLSKOLEM1 supplies only \(f=O(1)\) in the coarse sense of some polynomial exponent, not a certified numerical value of \(f\) or \(c_F\).

\section*{SAME\_REALIZATION and TC4 audit}
SAME\_REALIZATION remains load-bearing only if it requires more than semantic validity of the displayed predicate. If two different histories or source representations produce the same \((x,Q,\sigma,u)\) and any semantically valid \(z\) is accepted by \(R\), FULLSKOLEM1 permits RECOMP: construct another valid \(z\) rather than preserve the preferred history.

A surviving provenance or local-algebraic binding theorem must therefore establish a predicate stronger than semantic validity and prove that every relevant static SAT computation is forced to satisfy it. Merely naming the preferred source history does not create a computational obligation. This is the TC4 escape in reconstruction form.

\section*{PQ technique ledger}
Construction-specific collision extraction or special soundness remains \textbf{ANALOGY/CANDIDATE}: two conflicting valid openings could be useful only if they yield an unconditional forbidden algebraic witness in the exact SU/vMS realization. Generic cryptographic binding, random-oracle binding, and assumed collision resistance remain \textbf{FAILED TRANSFER} for an unconditional P-versus-NP separation. No reviewed PQ theorem overrides FULLSKOLEM1's semantic selector under \(P=NP\).

\section*{Relativization and algebrization audit}
FULLSKOLEM1 relativizes: for any oracle \(O\), if \(P^O=NP^O\), then \(PH^O=P^O\), and the same prefix searches are polynomial relative to \(O\). It therefore cannot be the missing nonrelativizing hinge. The argument also supplies no nonalgebrizing mechanism. Algebraic-oracle augmentation does not restore a semantic acquisition lower bound.

\section*{DUPLICATION / CONTRADICTION AUDIT}
FACTOR1 remains exact for fixed SAT-answer vectors. RAMCOMP1 remains a quantitative upper bound in its declared model. ENVELOPE1, RAMDOM1, CROSSDOM1, CERTREG1, TOTALENV1, and BOTTLENECK1 remain accounting results. FULLSKOLEM1 does not contradict them because it is qualitative under \(P=NP\), while those reports track explicit exponents for a fixed decider.

Generic quantifier commutation remains \textbf{FALSIFIED/RETIRED}. FULLSKOLEM1 never swaps \(\forall u\exists z\) into \(\exists z\forall u\); it constructs a uniform algorithm whose output may depend on the complete \(u\).

\section*{WHAT CHANGED RELATIVE TO ROOT MAIN}
Root main.tex currently records the displayed reconstruction prefix and says qualitative \(P=NP\Rightarrow PH=P\) is not a free quantitative compiler. That statement remains correct. The new point is complementary: although PH collapse gives no certified exponent, it is sufficient to rule out \emph{superpolynomial semantic acquisition hardness} for both \(\sigma\) and the \(u\mapsto z\) evaluator. Root main should therefore distinguish quantitative exponent accounting from qualitative semantic acquisition closure.

\section*{ADOPTED / REJECTED / RETIRED PRIOR COMPONENTS}
\textbf{Adopted:} FACTOR1; RAMCOMP1 at its quantitative upper-bound scope; ENVELOPE1 and descendants as accounting; STATICNP1 as an NP-staticization closure; SAME\_REALIZATION as a required source-matching discipline.

\textbf{Rejected:} promoting the size of an extensional Skolem table, the number of challenges \(u\), or the displayed alternation depth into a superpolynomial acquisition lower bound under \(P=NP\).

\textbf{FALSIFIED/RETIRED as generic hinges:} semantic Skolem-acquisition hardness for polynomially balanced fixed-level PH reconstruction; generic quantifier commutation; free canonical binding; TransportCert as a lower-bound theorem.

\section*{STALE OR OVERSTATED CLAIMS FOUND}
Any statement that the displayed \(\exists\sigma\forall u\exists z\forall y\) semantics alone may force superpolynomial acquisition under the \(P=NP\) contradiction is stale after FULLSKOLEM1. The quantitative warning that PH collapse does not provide a useful explicit exponent is not stale and must remain.

\section*{One next target}
Source-transcribe the strongest extra SAME\_REALIZATION requirement that is \emph{not} already represented inside \(R\). State it as a precise relation \(B(x,Q,\sigma,u,z,h)\) over source/history data \(h\). Then test:
\[
\text{(i) is }B\text{ semantically necessary for every static SAT decider?}
\]
\[
\text{(ii) does existence of a valid }(z,h)\text{ remain inside fixed-level PH?}
\]
If both answers are yes, FULLSKOLEM1 applies again after adjoining \(h\) as a polynomially balanced witness. A surviving route therefore needs either a necessity theorem for a genuinely stronger operation or a model boundary that cannot be absorbed into polynomially balanced PH search without changing the original SAT problem.

\section*{Strict status}
FULLSKOLEM1 is \textbf{PROVED/EXACT} as the stated conditional closure theorem. Application to a concrete vMS/SGL reconstruction is \textbf{CONDITIONAL} until the exact source predicate, witness lengths, and model are transcribed. SAME\_REALIZATION necessity, TC4 arbitrary-algorithm transfer, a source-matched lower bound, and genuinely nonrelativizing and nonalgebrizing hinges remain \textbf{OPEN}. P versus NP remains open.

\section*{Primary references}
E. Viola, On Probabilistic Time versus Alternating Time, ECCC TR05-137, 2005.

D. van Melkebeek and N. Mocelin Sdroievski, Instance-Wise Hardness and Refutation versus Derandomization for Arthur-Merlin Protocols, Computational Complexity 34, Article 15, 2025.

S. Akshay, O. Beyersdorff, S. Chakraborty, L. Kasche, M. Mahajan, and L. N. Spachmann, Proof Systems for QBF Synthesis: Extracting Skolem and Herbrand Functions, SAT 2026, LIPIcs 377, Article 3.
\end{document}
% END_OF_VERIFIED_REPORT::LSPV-3::FULLSKOLEM1